\chapter{Wie funktioniert Ausrüstung?}

\section{Was ist Ausrüstung?}

Ausrüstung sind Gegenstände, die ein Held bei sich trägt und verwenden kann. Sie wird durch Karten dargestellt.

\section{Ausrüstungsplätze}

Ein Held besitzt die Ausrüstungsplätze Kopf, zwei Hände, Körper, Gürtel und Füße. Nur Gegenstände, die an ihrem passenden Platz liegen, gelten als ausgerüstet und können verwendet werden.

Ein Held kann außerdem beliebig viel verstaute Ausrüstung mit sich tragen. Verstaute Ausrüstung gilt nicht als ausgerüstet und ihre Angaben und Effekte können nicht genutzt werden.

\subsection{Hände}

Ein Held hat zwei Hände. Manche Gegenstände benötigen eine Hand, andere beide Hände. Die ausgerüsteten Gegenstände dürfen zusammen höchstens zwei Hände belegen.

\section{Ausrüstung wechseln}

Mit einer \moveAction{} darf ein Held seine ausgerüsteten und verstauten Gegenstände beliebig neu anordnen. Er kann dabei mehrere Ausrüstungsplätze gleichzeitig verändern.

Der Körperplatz ist eine Ausnahme: Einen Gegenstand am Körperplatz auszurüsten, abzulegen oder gegen einen anderen zu tauschen kostet eine eigene \standardAction.

\section{Ausrüstungskarten lesen}

Waffen besitzen eine oder mehrere Angriffsoptionen. Jede Angriffsoption nennt die Angriffsfertigkeit, die Umrechnung von \star{} in Schadenswürfel und mögliche besondere Effekte.

Rüstungen und andere Gegenstände können eine Ausweichprobe (\evade) erlauben oder zusätzliche \blueDie{} dafür gewähren. Trägt ein Held keine Rüstung am Körperplatz, darf er mit einer Turnen-Probe ausweichen.

Schutz wird durch Schild-Symbole (\armor) dargestellt. Jedes \armor{} verhindert 1 Schaden. Rüstungsdurchdringung (\armorPierce) ignoriert die angegebene Anzahl an \armor{} des Ziels.

\section{Erschöpfte Karten}

Manche Effekte fordern den Spieler auf, eine Karte zu \textbf{erschöpfen (\exhaust)}. Dazu wird die Karte umgedreht.

Eine erschöpfte Karte bleibt ausgerüstet und belegt weiterhin ihren Ausrüstungsplatz. Da ihre Vorderseite nicht sichtbar ist, können jedoch keine ihrer Angaben, Werte oder Effekte genutzt werden.

Erschöpfte Karten werden spätestens bei der nächsten Rast wieder bereitgemacht. Andere Regeln können sie früher bereitmachen.

\section{Weiterführende Informationen}

Im Anhang befinden sich die vollständige Ausrüstungsliste und die Symbolübersicht.
